\begin{center}
    {\LARGE \textbf{2015分析化学}} \\[2ex]
    {\large 时量180分钟 \quad 满分90分}
\end{center}

\vspace{0.5cm}
\noindent\textbf{注意事项:}
\begin{enumerate}[label=\arabic*., parsep=0pt, itemsep=0pt]
    \item 答题前,考生先将自己的姓名、学号、考场号、座位号填写在答题卡上,将条形码准确粘贴在考生信息条形码粘贴区。
    \item 考试结束后,将本试卷与答题纸一并收回。
    \item 回答各个问题时,将答案写在答题纸上,写在本试卷上无效。
\end{enumerate}

\vspace{0.5cm}
\noindent\textbf{一、简答题(42 pts.)}

\begin{enumerate}[label=\arabic*., itemsep=1.5ex]
    \item 写出用$\ce{K2Cr2O7}$标准溶液标定$\ce{Na2S2O3}$的反应方程式
    \begin{enumerate}[label=(\arabic*), noitemsep, topsep=0pt, partopsep=0pt]
        \item 
        \item 
    \end{enumerate}
    \item 写出下列溶液的质子条件式：
    \begin{enumerate}[label=(\arabic*), noitemsep, topsep=0pt, partopsep=0pt]
        \item $0.1\ \mathrm{mol/L}\ \ce{NH4Ac}$溶液
        \item $0.1\ \mathrm{mol/L}\ \ce{H2SO4}$溶液
    \end{enumerate}
    \item 某物质的水溶液$100\ \mathrm{mL}$，用$5$份$10\ \mathrm{mL}$萃取剂溶液连续萃取$5$次，总萃取率为$87\%$，则该物质在此萃取体系中的分配比是多少？
    \item 某同学配制$0.02\ \mathrm{mol/L}\ \ce{Na2S2O3}\ 500\ \mathrm{mL}$，方法如下：
    
    在分析天平上准确称取$\ce{Na2S2O3 * 5H2O}\ 2.482\ \mathrm{g}$，溶于蒸馏水中，加热煮沸，冷却，转移至$500\ \mathrm{mL}$容量瓶中，加蒸馏水定容摇匀，保存待用。请指出其错误。
    \item 已知半反应$\ce{VO2^+ + 2H+ + e^- -> VO^{2+} + 2H2O}$的标准电极电位$E^\ominus=1.00\ \mathrm{V}$，计算在硫酸介质中，$\ce{H+}$浓度为$0.1\ \mathrm{mol/L}$的溶液中电对$\ce{VO2^+/VO^{2+}}$的条件电位$E^{\ominus}$.
    \item 用纸色谱法分离试液中的$\ce{Co^{2+}}$，$\ce{Ni^{2+}}$，已知$\ce{Co^{2+}}$的$R_f=0.49$，$\ce{Ni^{2+}}$的$R_f=0.01$，欲使它们斑点中心间的距离为$3\ \mathrm{cm}$，溶剂前沿与顶端相距最少$1\ \mathrm{cm}$，问滤纸条最短应多长。
    \item 测定某金属中硫的质量分数$w(\mathrm{S}\%)$，结果为$0.097$，$0.098$，$0.100$，$0.102$，$0.098$，$0.099$。计算置信度为$90\%$时平均值的置信区间。(已知$n=6$，置信度为$90\%$时$t=2.015$)
\end{enumerate}

\vspace{0.5cm}
\noindent\textbf{二、计算题(48 pts. total)}

\begin{enumerate}[label=\arabic*., itemsep=1.5ex]
    \item (12 pts.) 计算在$\mathrm{pH}=1.70$的盐酸溶液中$\ce{CaF2}$的溶解度。(已知$\ce{CaF2}$的$K_{\mathrm{sp}}=3.2\times10^{-11}$，$\ce{HF}$的$K_{\mathrm{a}}=6.6\times10^{-4}$)
    \item (12 pts.) 求反应$\ce{Fe^{3+} + Ti^{3+} = Fe^{2+} + Ti^{4+}}$的平衡常数。当用$0.0100\ \mathrm{mol/L}\ \ce{TiCl3}$滴定$0.0100\ \mathrm{mol/L}\ \ce{Fe^{3+}}$溶液，直到$\ce{KSCN}$不与试液显现红色，此时$[\ce{Fe^{3+}}]=10^{-5}\ \mathrm{mol/L}$，求溶液中$[\ce{Ti^{4+}}]/[\ce{Ti^{3+}}]$的比值。
    \\
    （已知$E^{\ominus}(\ce{Fe^{3+}/Fe^{2+}})=0.77\ \mathrm{V}$，$E^{\ominus}(\ce{TiO^{2+}/Ti^{3+}})=0.1\ \mathrm{V}$）
    \item (12 pts.) 用$0.1000\ \mathrm{mol/L}\ \ce{HCl}$溶液滴定$20.00\ \mathrm{mL}\ 0.1000\ \mathrm{mol/L}$羟胺，计算化学计量点及滴定突跃的$\mathrm{pH}$值。(已知羟胺的$\mathrm{p}K_{\mathrm{b}}=8.04$)
    \item (12 pts.) 浓度为$0.02\ \mathrm{mol/L}$的$\ce{Cu^{2+}}$溶液，其$\mathrm{pH}=10.0$，$c_{\ce{NH3}}=0.02\ \mathrm{mol/L}$，用$\mathrm{PAN}$为指示剂，用等浓度的$\mathrm{EDTA}$滴定$\ce{Cu^{2+}}$，终点时$\mathrm{pCu_{\mathrm{ep}}}=11.2$，计算终点误差。
    \\
    (已知$\lg K_{\ce{CuY}}=18.8$，$\mathrm{p}K_{\mathrm{b}}(\ce{NH3})=4.74$，$\mathrm{pH}=10.0$时，$\lg\alpha_{\mathrm{Y(H)}}=0.5$，$\lg\alpha_{\mathrm{Cu(OH)}}=0.8$；$\ce{Cu^{2+}-NH3}$络合物的$\lg\beta_1-\lg\beta_4$分别为$4.13$，$7.61$，$10.48$，$12.59$)
\end{enumerate}